\section{Preliminaries}\label{sec: preliminaries}

\subsection{Simple matroids and geometric lattices}

In order to avoid confusion, we do not use geometrical terminology such as ``point" and ``hyperplane" for matroids and we call an element of a matroid an \textit{atom} and a flat of corank $ 1 $ of a matroid a \textit{coatom}.
This is lattice theoretic terminology due to the following well-known theorem.

\begin{Theorem}[see {\cite[p.~54, Theorem 2]{welsh2010matroid}} for example]\label{simple matroid geometric lattice correspondence}
The correspondence between a simple matroid and its lattice of flats is a bijection between simple matroids and geometric lattices.
\end{Theorem}

Thus any properties about simple matroids are translated into properties of geometric lattices, and vice versa.
For example, the contraction and the restriction of matroids are just intervals in the lattice of flats as follows.

\begin{Proposition}[{\cite[Proposition 3.3.8]{oxley2011matroid}}]\label{Oxley interval}
Let $ X $ be a flat of a matroid $ M $. Then
\begin{enumerate}[$(1)$]\itemsep=0pt
\item\label{Oxley interval 1} $ L(M/X) \simeq \big[X, \hat{1}\big] = \{ F \in L(M) \,|\, X \leq F \} $.
\item\label{Oxley interval 2} $ L(M|X) \simeq \big[\hat{0}, X\big] = \{ F \in L(M) \,|\, F \leq X \} $.
\end{enumerate}
\end{Proposition}

Note that the contraction of a simple matroid is not simple in general.
However, we can associate the simple matroid $ \si(M) $ with a matroid $ M $ and this operation does not affect the lattice of flats, that is, $ L(M) \simeq L(\si(M)) $.

If $ \mathcal{A} $ be an arrangement and $ M(\mathcal{A}) $ the simple matroid on $ \mathcal{A} $, then for any hyperplane $ H \in \mathcal{A} $, we have $ M\big(\mathcal{A}^{H}\big) \simeq \si(M(\mathcal{A})/H) $, where $ \mathcal{A}^{H} $ denotes the \textit{restriction} defined by
\begin{gather*}
\mathcal{A}^{H} \coloneqq \{ K \cap H \,|\, K \in \mathcal{A}\setminus\{H\}\}.
\end{gather*}

Note that the restriction of an arrangement does \emph{not} correspond to the restriction of a~matroid but the simplification of the contraction.
The restriction of a matroid corresponds to the localization of an arrangement.

\subsection{Characteristic polynomials}

Let $ L $ be a geometric lattice.
The \textit{characteristic polynomial} $ \chi(L,t) \in \mathbb{Z}[t] $ is defined by
\begin{gather*}
\chi(L,t) \coloneqq \sum_{X \in L}\mu(X)t^{r(\hat{1}) - r(X)},
\end{gather*}
where $ r $ denotes the rank function of $ L $ and $ \mu \colon L \to \mathbb{Z} $ denotes the one-variable M\"{o}bius function of $ L $ defined recursively by{\samepage
\begin{gather*}
\mu(X) \coloneqq \begin{cases}
 1 & \text{if } X = \hat{0}, \vspace{1mm}\\
-\displaystyle\sum_{Y < X}\mu(Y) & \text{otherwise.}
\end{cases}
\end{gather*}
The \textit{characteristic polynomial} of $ M $ is defined by $ \chi(M,t) \coloneqq \chi(L(M),t) $.}

The \textit{intersection lattice} $ L(\mathcal{A}) $ of a central arrangement $ \mathcal{A} $ is defined by
\begin{gather*}
L(\mathcal{A}) \coloneqq \bigg\{ \bigcap_{H \in \mathcal{B}}H \,\bigg|\, \mathcal{B} \subseteq \mathcal{A} \bigg\}
\end{gather*}
with a partial order by reverse inclusion.
Note that $ L(\mathcal{A}) $ is a geometric lattice and naturally isomorphic to $ L(M(\mathcal{A})) $.
The \textit{characteristic polynomial} $ \chi(\mathcal{A}, t) \in \mathbb{Z}[t] $ is defined by
\begin{gather*}
\chi(\mathcal{A},t) \coloneqq \sum_{X \in L(\mathcal{A})}\mu(X)t^{\dim X}.
\end{gather*}
Since the rank function of $ L(\mathcal{A}) $ is given by the codimension, $ \chi(\mathcal{A},t) = t^{\ell-r}\chi(M(\mathcal{A}), t) $, where $ \ell $ is the dimension of the ambient space and $ r $ is the rank of $ L(\mathcal{A}) $.
When $ \ell = r $ we say that $ \mathcal{A} $ is \textit{essential}.
It is well known that there exists an essential arrangement $ \mathcal{A}_{0} $ for every arrangement~$ \mathcal{A} $ such that $ L(\mathcal{A}) = L(\mathcal{A}_{0}) $ and $ \mathcal{A} $ is free if and only if~$ \mathcal{A}_{0} $ is free (see \cite{orlik1992arrangements} for details).

\subsection{Modularity}

We excerpt some conditions equivalent to modularity from Brylawski \cite{brylawski1975modular-totams}.
\begin{Proposition}[Brylawski {\cite[Theorem~3.3]{brylawski1975modular-totams}}]\label{brylawski modular equivalent}
Let $ X $ be an element of a geometric lattice~$ L $.
Then the following conditions are equivalent.
\begin{enumerate}[$(1)$]\itemsep=0pt
\item $ X $ is modular.
\item\label{brylawski modular equivalent 2} For $ Y \leq Z $ in $ L $, $ Y \vee (X \wedge Z) = (Y \vee X) \wedge Z $.
\item\label{brylawski modular equivalent 3} For all $ Y \in L $, $ [X \wedge Y, X] \simeq [Y, X \vee Y] $.
\item\label{brylawski modular equivalent 4} For any atom $ e \not\leq X $, $ [\hat{0}, X] \simeq [e, X \vee e] $ and $ X \vee e $ is modular in $ \big[e, \hat{1}\big] $.
\end{enumerate}
\end{Proposition}


\begin{Theorem}[Brylawski {\cite[Theorem 3.11]{brylawski1975modular-totams}} (the modular short-circuit axiom)]\label{Brylawsky modular short-circuit axiom}
Let $ M $ be a~simple matroid on the ground set $ E $ and $ X \subseteq E $ a nonempty subset.
Then $ X $ is a modular flat of $ M $ if and only if for every circuit $ C $ of $ M $ and an atom $ e \in C \setminus X $ there exist an atom $ x \in X $ and a circuit $ C^{\prime} $ such that $ e \in C^{\prime} \subseteq \{x\} \cup (C \setminus X) $.
\end{Theorem}

\begin{Theorem}[{Brylawski \cite[Corollary 3.4]{brylawski1975modular-totams}}]\label{Brylawski modular coatom}
A coatom $ X $ of a simple matroid on $ E $ is modular if and only if for any distinct two atoms $ e, e^{\prime} \in E\setminus X $ there exists $ e^{\prime\prime} \in X $ such that $ \{e,e^{\prime},e^{\prime\prime}\} $ forms a circuit.
\end{Theorem}

We give some properties of modularity required in this article.
\begin{Proposition}[{Jambu--Papadima \cite[Lemmas~1.3 and~1.9]{jambu1998generalization-t}}]\label{Jambu-Papadima}
Let $ X $ be a modular coatom of a simple matroid $ M $ on $ E $. Then for any two distinct atoms $a,b \in E\setminus X $, there exists unique $ f(a,b) \in X $ such that $ a,b,f(a,b) $ form a circuit.
Moreover, for any three distinct atoms $ a,b,c \in E\setminus X $, the atoms $ f(a,b)$, $f(a,c)$, $f(b,c) $ form a circuit.
\end{Proposition}

\begin{Proposition}[Brylawski {\cite[Proposition 3.5]{brylawski1975modular-totams}}]\label{Brylawski modular tower}
Let $ X $ be a modular flat of a simple matroid~$ M $ and~$ Y $ a modular flat of the restriction~$ M|X $, then~$ Y $ is a modular flat of~$ M $.
\end{Proposition}

\begin{Proposition}[Brylawski {\cite[Proposition~3.6]{brylawski1975modular-totams}}]\label{Brylawski modular intersection}
Let $ X $ and $ Y $ be modular flats of a simple matroid.
Then $ X \cap Y $ is a modular flat.
\end{Proposition}

\begin{Proposition}[Probert {\cite[Corollary 4.2.8]{probert2018chordality}}]\label{Probert}
Every modular flat of a round matroid is round.
\end{Proposition}

\begin{Theorem}[Stanley {\cite[Theorem 2]{stanley1971modular-au}}]\label{stanley division}
If $ X $ is a modular element of a geometric lattice $ L $, then $ \chi\big(\big[\hat{0},X\big],t\big) $ divides $ \chi(L,t) $.
\end{Theorem}

\subsection{Divisionality}

The following theorem plays an important role in this article.
\begin{Theorem}[Abe {\cite[Theorem 1.1 (division theorem)]{abe2016divisionally-im}}]\label{Abe division}
An arrangement $ \mathcal{A} $ is free if there exists $ H \in \mathcal{A} $ such that $ \mathcal{A}^{H} $ is free and $ \chi\big(\mathcal{A}^{H},t\big) $ divides $ \chi(\mathcal{A},t) $.
\end{Theorem}

Theorem \ref{Abe division} leads to the concepts of a divisional flag, which is originally defined for arrangements.
However, we define it for simple matroids as follows.

\begin{Definition}\label{definition divisional flag}
A \textit{divisional flag} of a simple matroid $ M $ of rank $ n $ is a sequence of flats \begin{align*}
\varnothing = X_{0} \subseteq X_{1} \subseteq \dots \subseteq X_{n} = E
\end{align*}
such that $ r(X_{i}) = i $ for $ i \in \{0, \dots, n \} $ and $ \chi(M/X_{i+1},t) \,|\, \chi(M/X_{i}, t) $ for $ i \in \{0, \dots, n-1 \} $.
\end{Definition}
Note that since $ \chi(M/X_{n},t)=1 $ and $ \chi(M/X_{n-1},t) = t-1 $ in the definition, the conditions $ \chi(M/X_{n},t) \,|\, \chi(M/X_{n-1},t) $ and $ \chi(M/X_{n-1},t) \,|\, \chi(M/X_{n-2},t) $ are always satisfied.

\begin{Definition}\label{definition divisional freeness}
An arrangement called \textit{divisionally free} if the matroid $ M(\mathcal{A}) $ has a divisional flag.
\end{Definition}
Note that, by Theorem \ref{Abe division}, every divisionally free arrangement is free. In order to find a~divisional flag, define a divisional atom as follows.

\begin{Definition}An atom $ e $ of a simple matroid $ M $ is called \textit{divisional} if $ \chi(M/e,t) $ divides $ \chi(M,t) $.
\end{Definition}

\begin{Proposition}\label{divisional flag}
A nonempty simple matroid $ M $ has a divisional flag if and only if there exists a divisional atom $ e $ such that $ \si(M/e) $ has a divisional flag.
\end{Proposition}
\begin{proof}
Suppose that $ M $ has a divisional flag $ \varnothing = X_{0} \subseteq X_{1} \subseteq \dots \subseteq X_{n} = E $.
Then $ e \coloneqq X_{1} $ is a divisional atom and the images of $ X_{1}, \dots, X_{n} $ under the isomorphism $ [e, \hat{1}] \simeq L(\si(M/e)) $ given by Proposition \ref{Oxley interval}(\ref{Oxley interval 1}) form a divisional flag of $ \si(M/e) $.
The converse holds by a similar argument.
\end{proof}

Abe \cite[Proposition 5]{abe2017restrictions-pilt} proved that every supersolvable arrangement is divisionally free by constructing a divisional flag from a saturated chain of modular flats. The following lemma is a~generalization for simple matroids.

\begin{Lemma}\label{modular coatom divisional}
Let $ M $ be a simple matroid on the ground set $ E $ and $ X $ a modular coatom of $ M $.
Then every atom $ e \in E \setminus X $ is divisional and $ \si(M/e) \simeq M|X $.
In particular, every supersolvable matroid has a divisional flag.
\end{Lemma}
\begin{proof}
Take an atom $ e \in E\setminus X $, and $ e $ is a complement of $ X $, that is, $ X \wedge e = \hat{0} $ and $ X \vee e = \hat{1} $.
By Proposition \ref{brylawski modular equivalent}(\ref{brylawski modular equivalent 3}), we have $ \big[\hat{0}, X\big] \simeq \big[e, \hat{1}\big] $, which implies $ M|X \simeq \si(M/e) $.
Moreover, by Theo\-rem~\ref{stanley division}, the characteristic polynomial $ \chi(M/e,t) = \chi(\si(M/e),t) = \chi(M|X,t) = \chi\big(\big[\hat{0}, X\big], t\big) $ divides $ \chi(M,t) $ and hence $ e $ is divisional.

When $ M $ is supersolvable, $ \si(M/e) $ is supersolvable.
Therefore $ M $ has a divisional flag by induction and Proposition \ref{divisional flag}.
\end{proof}

\subsection{Modular joins}

We review some properties of modular joins and will show a relation between modular joins and divisional atoms.

\begin{Proposition}[see also Ziegler {\cite[Lemma 3.10]{ziegler1991binary-potams}}]\label{Ziegler case of modular coatom}
Let $ X $ be a minimal flat of a simple matroid $ M $ such that $ M $ is a modular join $ M = P_{X}(M_{1}, M_{2}) $ over $ X $.
If $ M_{1} $ has a modular coatom, then $ M_{1} $ has a divisional atom not belonging to $ X $.
\end{Proposition}
\begin{proof}
Let $ Z \subseteq E_{1} $ be a modular coatom of $ M_{1} $, where $ E_{1} $ denotes the ground set of $ M_{1} $.
By Lemma \ref{modular coatom divisional}, every element in $ E_{1}\setminus Z $ is a divisional atom of $ M_{1} $.
Assume that $ E_{1}\setminus Z \subseteq X $.
Then $ Z \cap X \subsetneq X $ and $ M $ is a modular join $ M = P_{Z \cap X}(M|Z, M_{2}) $ over $ Z \cap X $, which is a contradiction to the minimality of $ X $.
Hence $ E_{1}\setminus Z \not\subseteq X $ and every element $ E_{1}\setminus(Z \cup X) $ is a desired atom.
\end{proof}

\begin{Theorem}[Brylawski {\cite[Theorem 7.8]{brylawski1975modular-totams}}]\label{Brylawski decomposition}
Let $ M = P_{X}(M_{1}, M_{2}) $ be a modular join.
Then
\begin{align*}
\chi(M,t) = \frac{\chi(M_{1},t) \, \chi(M_{2},t)}{\chi(M|X,t)}.
\end{align*}
\end{Theorem}

The following proposition is essentially due to Brylawski for generalized parallel connections.
However, since we treat the special case of modular joins, we give a proof of the proposition below.
\begin{Proposition}[Brylawski {\cite[Theorem 5.11.4]{brylawski1975modular-totams}}]\label{Brylawski contraction}
Let $ M $ be a modular join $ M=P_{X}(M_{1},M_{2}) $ and $ e $ an atom of $ M_{1} $ not belonging to $ X $.
Then $ \si(M/e) $ is isomorphic to a modular join of $ \si(M_{1}/e) $ and $ M_{2}, $ and
\begin{align*}
\chi(\si(M/e),t) = \dfrac{\chi(\si(M_{1}/e),t) \chi(M_{2}, t)}{\chi(M|X, t)}.
\end{align*}
\end{Proposition}
\begin{proof}
Let $ E_{1} $ and $ E_{2} $ be the ground sets of $ M_{1} $ and $ M_{2} $, which are modular flats of~$ M $.
Take an atom $ e \in E_{1} \setminus X $.
The matroid $ \si(M/e) $ corresponds the interval $ \big[e, \hat{1}\big] $ of~$ L(M) $ under the correspondence mentioned in Proposition \ref{simple matroid geometric lattice correspondence}.
Note that $ E_{1} $ is modular in $ \big[e, \hat{1}\big] $ by Proposition~\ref{brylawski modular equivalent}(\ref{brylawski modular equivalent 2}) and $ E_{2} \vee e $ is modular in $ \big[e, \hat{1}\big] $ by Proposition~\ref{brylawski modular equivalent}(\ref{brylawski modular equivalent 4}).

The atoms of $ \si(M/e) $ are identified with the atoms of the interval $ \big[e, \hat{1}\big] $.
These atoms coincide with $ \{ e \vee e^{\prime} \,|\, e^{\prime} \in E\setminus \{e\} \} $, where $ E = E_{1} \cup E_{2} $ denotes the ground set of $ M $.
If $ e^{\prime} \in E_{1} $, then $ e \vee e^{\prime} \leq E_{1} $.
Suppose that $ e^{\prime} \in E_{2} $.
Then $ e \vee e^{\prime} \leq e \vee E_{2} $.
Thus $ \si(M/e) $ is a modular join of matroids corresponding to $ [e, E_{1}] $ and $ [e, E_{2} \vee e] $.
The matroid corresponding to $ [e, E_{1}] $ is isomorphic to $ \si(M_{1}/e) $.
By Proposition~\ref{brylawski modular equivalent}(\ref{brylawski modular equivalent 3}), $ [e, E_{2} \vee e] \simeq \big[\hat{0}, E_{2}\big] $. Hence the matroid corresponding to $ [e, E_{2} \vee e] $ is isomorphic to~$ M_{2} $.
Thus $ \si(M/e) $ is isomorphic to a modular join of $ \si(M_{1}/e) $ and~$ M_{2} $.

By Proposition \ref{brylawski modular equivalent}(\ref{brylawski modular equivalent 2}), $ E_{1} \wedge (E_{2} \vee e) = (E_{1} \wedge E_{2}) \vee e = X \vee e $.
Using Theorem \ref{Brylawski decomposition} and Proposition \ref{brylawski modular equivalent}(\ref{brylawski modular equivalent 3}), we have
\begin{align*}
\chi(\si(M/e),t)& = \dfrac{\chi([e, E_{1}],t) \, \chi([e, E_{2}\vee e], t)}{\chi([e,X \vee e]), t}
= \dfrac{\chi([e, E_{1}],t) \chi\big(\big[\hat{0}, E_{2}\big], t\big)}{\chi\big(\big[\hat{0},X\big]\big), t} \\
&= \dfrac{\chi(\si(M_{1}/e),t) \chi(M_{2}, t)}{\chi(M|X, t)}.\tag*{\qed}
\end{align*}\renewcommand{\qed}{}
\end{proof}

\begin{Lemma}\label{divisional atom of modular join}
Let $ M $ be a modular join $ M = P_{X}(M_{1},M_{2}) $.
Every divisional atom of $ M_{1} $ not belonging to $ X $ is a divisional atom of $ M $.
\end{Lemma}
\begin{proof}
Let $ e $ be a divisional atom of $ M_{1} $ such that $ e \not\in X $.
Then there exists an integer $ a $ such that $ \chi(M_{1},t) = (t-a)\chi(\si(M_{1}/e),t) $.
Using Proposition \ref{Brylawski contraction}, we have
\begin{align*}
\chi(M,t) = \dfrac{\chi(M_{1},t) \, \chi(M_{2},t)}{\chi(M|X,t)}
= \dfrac{(t-a) \, \chi(\si(M_{1}/e),t) \, \chi(M_{2},t)}{\chi(M|X,t)}
= (t-a)\chi(\si(M/e),t).
\end{align*}
Thus $ e $ is a divisional atom of $ M $.
\end{proof}

\section{Proof of main theorems}\label{sec: proof}

\subsection{Proof of Theorem \ref{main divisionally free}}
\begin{Lemma}\label{restriction modular flat}
The class $ \mathcal{ME} $ is closed under taking restrictions to modular flats.
\end{Lemma}
\begin{proof}
Let $ M \in \mathcal{ME} $ and $ X $ a modular flat of $ M $.
We proceed by induction on the rank of $ M $.
The case $ r(M)=0 $ is trivial.
Hence we suppose that $ r(M) \geq 1 $.

First assume that $ M $ has a modular coatom $ Z $ such that $ M|Z \in \mathcal{ME} $.
If $ X \subseteq Z $, then $ X $ is a~modular flat of $ M|Z $.
By the induction hypothesis, $ M|X = (M|Z)|X \in \mathcal{ME} $.
Assume $ X \not\subseteq Z $.
Then $ X \vee Z = E $, the ground set of $ M $.
By Proposition~\ref{Brylawski modular intersection}, $ X \cap Z $ is a modular flat of $ M $ and hence $ M|Z $.
By the induction hypothesis, $ M|(X \cap Z) \in \mathcal{ME} $.
Moreover, by the modularity,
\begin{gather*}
r(X) - r(X \cap Z) = r(X \vee Z) - r(Z) = r(E) - r(Z) = 1.
\end{gather*}
Therefore $ X \cap Z $ is a modular coatom of $ M|X $ and hence $ M|X \in \mathcal{ME} $.

Next we suppose that $ M $ is a modular join $ M = P_{Y}(E_{1}, E_{2}) $ over a round flat $ Y $ with $ M|E_{i} \in \mathcal{ME} $ for $ i = 1,2 $.
If $ X \subseteq E_{i} $ for some $ i $, then $ M|X = (M|E_{i})|X \in \mathcal{ME} $ by the induction hypothesis.
Otherwise, $ X_{i} \coloneqq X \cap E_{i} \neq \varnothing $ is a proper subset of $ X $ and $ M|X_{i} = (M|E_{i})|X_{i} \in \mathcal{ME} $ by the induction hypothesis for $ i = 1,2 $.
Since both $ X_{1} $ and $ X_{2} $ are modular by Proposition~\ref{Brylawski modular intersection} and $ X = X_{1} \cup X_{2} $, it follows that $ M|X $ is a modular join of $ M|X_{1} $ and $ M|X_{2} $.
Moreover $ X_{1} \cap X_{2} = X \cap Y $ is round by Proposition~\ref{Probert}.
Therefore $ M|X \in \mathcal{ME} $.
\end{proof}

\begin{Theorem}\label{main matroid}
Every nonempty simple matroid $ M \in \mathcal{ME} $ has a divisional atom $ e $ such that $ \si(M/e) \in \mathcal{ME} $.
\end{Theorem}
\begin{proof}
We will proof the following claims by induction on the rank of $ M $.
\begin{enumerate}[$(i)$]\itemsep=0pt
\item\label{main matroid claim i} If $ M $ has a modular coatom, then there exists a divisional atom $ e $ such that $ \si(M/e) \in \mathcal{ME} $.
\item\label{main matroid claim ii} If $ X $ is a minimal round flat of $ M $ such that $ M $ is a modular join $ M = P_{X}(E_{1}, E_{2}) $.
Then, for each $ i=1,2 $, there exists a divisional atom $ e_{i} \in E_{i}\setminus X $ such that $ \si(M/e_{i}) \in \mathcal{ME} $.
\end{enumerate}

First suppose that $ r(M) = 1 $, that is, the ground set of $ M $ is a singleton.
Then only the case~(\ref{main matroid claim i}) occurs and the atom of $ M $ satisfies the assertion.

Now suppose that $ r(M) \geq 2 $.
If $ M $ has a modular coatom $ X $, then every atom $ e \in E \setminus X $ is divisional and $ \si(M/e) \simeq M|X \in \mathcal{ME} $ by Lemmas~\ref{modular coatom divisional} and~\ref{restriction modular flat}. Thus the assertion holds.

Next we suppose that $ M $ is a modular join.
We assume that $ X $ is a minimal round flat of $ M $ such that $ M = P_{X}(M_{1},M_{2}) $.
Since every modular flat in $ X $~is also round by Proposition~\ref{Probert}, $ X $~is a minimal flat such that $ M $ is a modular join over $ X $.

We will show that $ M_{1} $ has a divisional atom $ e_{1} $ not belonging to $ X $ such that $ \si(M_{1}/e_{1}) \in \mathcal{ME} $.
Note that $ M_{1} $ is a member of $ \mathcal{ME} $ by Lemma~\ref{restriction modular flat}.
Assume that $ M_{1} $ has a modular coatom.
Then, by Proposition~\ref{Ziegler case of modular coatom}, $ M_{1} $ has a divisional atom not belonging to $ X $ such that $ \si(M_{1}/e_{1}) \in \mathcal{ME} $.
Hence we may assume that $ M_{1} $ has a minimal round flat $ Y $ such that $ M_{1} $ is a modular join $ M_{1} = P_{Y}(F,F^{\prime})$.
Since $ X $ is round, we have $ F \supseteq X $ or $ F^{\prime} \supseteq X $.
Without loss of generality, we may assume that $ F^{\prime} \supseteq X $.
By the induction hypothesis, $ M_{1} $ has a divisional atom $ e_{1} \in F\setminus Y $ such that $ \si(M_{1}/e_{1}) \in \mathcal{ME} $.
Assume that $ e_{1} \in X $.
Then $ e_{1} \in F^{\prime} $ and hence $ e_{1} \in F \cap F^{\prime} = Y $, which contradicts $ e_{1} \not\in Y $.
Thus $ M_{1} $ has a divisional atom $ e_{1} $ not belonging to $ X $ such that $ \si(M_{1}/e_{1}) \in \mathcal{ME} $.

By Lemma~\ref{divisional atom of modular join}, $ e_{1} $ is a divisional atom of $ M $.
Moreover, by Proposition \ref{Brylawski contraction}, $ \si(M/e_{1}) $ is isomorphic to a modular join of $ \si(M_{1}/e_{1}) $ and $ M_{2} $, and hence $ \si(M/e_{1}) \in \mathcal{ME} $.
\end{proof}

\begin{proof}[Proof of Theorem~\ref{main divisionally free}] Use Theorem \ref{main matroid} and Proposition \ref{divisional flag}.
\end{proof}
